\documentclass[10pt,twocolumn,letterpaper]{article}

\usepackage[spanish,es-nodecimaldot,es-noshorthands,es-tabla]{babel}
\usepackage[margin=2cm,columnsep=0.8cm]{geometry}
\usepackage{amsmath,amssymb,amsthm}
\usepackage{graphicx}
\usepackage{booktabs}
\usepackage{array}
\usepackage{enumitem}
\usepackage{algorithm}
\usepackage[noend]{algpseudocode}
\usepackage{pgfplots}
\usepackage{xcolor}
\usepackage{caption}
\usepackage{hyperref}

\pgfplotsset{compat=1.18}
\captionsetup{font=small,labelfont=bf}
\setlist{itemsep=2pt,topsep=3pt}
\hypersetup{colorlinks=true,linkcolor=black,citecolor=black,urlcolor=black}

\newtheorem{definicion}{Definición}
\newtheorem{lema}{Lema}
\newtheorem{teorema}{Teorema}
\newtheorem{proposicion}{Proposición}
\newtheorem{observacion}{Observación}

\floatname{algorithm}{Algoritmo}
\algrenewcommand\algorithmicfor{\textbf{para}}
\algrenewcommand\algorithmicdo{\textbf{hacer}}
\algrenewcommand\algorithmicwhile{\textbf{mientras}}
\algrenewcommand\algorithmicif{\textbf{si}}
\algrenewcommand\algorithmicthen{\textbf{entonces}}
\algrenewcommand\algorithmicreturn{\textbf{retornar}}
\algrenewcommand\algorithmicrequire{\textbf{Entrada:}}
\algrenewcommand\algorithmicensure{\textbf{Salida:}}

\newcommand{\costo}{\operatorname{costo}}
\newcommand{\N}{\mathbb{N}}

\begin{document}

\twocolumn[
\begin{center}
  {\includegraphics[height=2.4cm]{logo-univalle.png}}\par\vspace{0.4cm}
  {\LARGE Primer proyecto -- ADA II}\par\vspace{0.15cm}
  {\Large Programación de tareas con penalización por retraso}\par\vspace{0.1cm}
  {\large Programación dinámica y algoritmos voraces para $1\,\|\,\sum w_jT_j$}\par\vspace{0.4cm}
  {\large Pablo Nicolás Marín González \,(2459440) \qquad Kevin Steven Rico González \,(2411024)}\par\vspace{0.15cm}
  Análisis y Diseño de Algoritmos II \quad·\quad Universidad del Valle \quad·\quad 2026-2\par
  \vspace{0.6cm}
\end{center}
]

\begin{abstract}
Se estudia el problema de ordenar $n$ tareas en una sola máquina para minimizar
la suma ponderada de sus retrasos, $\sum w_jT_j$. Se demuestra que su versión de
decisión pertenece a NP y se discute su NP-dificultad en sentido fuerte. Se
propone una solución exacta por programación dinámica sobre subconjuntos, de
complejidad $O(2^n n)$, y tres algoritmos voraces (EDD, WSPT y WMDD) de
complejidad $O(n\log n)$ u $O(n^2)$. Ambos enfoques se implementan en C++ y se
comparan sobre 315 instancias aleatorias con $n$ entre 4 y 24. La programación
dinámica coincide con la fuerza bruta en todos los casos verificables y su
tiempo se duplica con cada tarea adicional; el voraz WMDD obtiene brechas
medianas de 0.2\,\% a 1.2\,\% cuando las fechas límite son exigentes, pero
ningún voraz garantiza el óptimo.
\end{abstract}

%=======================================================================
\section{Análisis teórico del problema}
%=======================================================================

\subsection{Descripción informal}

Una sola máquina (o persona) debe ejecutar $n$ tareas, una a la vez y sin
interrumpirlas. Cada tarea tarda un tiempo, tiene una fecha en la que debería
estar terminada y una penalización por cada unidad de tiempo que se entregue
tarde. La pregunta es en qué orden ejecutarlas para que la penalización total
sea mínima.

\subsection{Modelo formal}

\begin{definicion}[Instancia]
Una instancia es un conjunto de tareas $J=\{1,\dots,n\}$ donde cada tarea $j$
tiene una \emph{duración} $p_j\in\N_{>0}$, una \emph{fecha límite}
$d_j\in\N$ y un \emph{peso} o penalización por unidad de retraso
$w_j\in\N$. Todas las tareas están disponibles en el instante $0$.
\end{definicion}

\begin{definicion}[Programación y costo]
Una programación es una permutación $\pi=(\pi_1,\dots,\pi_n)$ de $J$. La tarea
en la posición $k$ termina en
\[
C_{\pi_k} = \sum_{i=1}^{k} p_{\pi_i},
\]
su \emph{retraso} es $T_j=\max(0,\;C_j-d_j)$ y el costo de $\pi$ es
\[
\costo(\pi)=\sum_{j=1}^{n} w_j\,T_j=\sum_{j=1}^{n} w_j\max(0,\;C_j-d_j).
\]
\end{definicion}

\paragraph{Problema de optimización.} Encontrar
$\pi^\ast=\arg\min_{\pi\in S_n}\costo(\pi)$, donde $S_n$ es el conjunto de las
$n!$ permutaciones. En la notación de Graham se escribe
$1\,\|\,\sum w_jT_j$: una máquina ($1$), sin restricciones adicionales
(campo vacío) y objetivo de retraso ponderado total.

No es necesario considerar tiempos ociosos ni interrupciones: dejar la máquina
quieta solo retrasa las tareas siguientes, y como todas están disponibles desde
$t=0$, toda programación óptima puede tomarse sin huecos. Así, el espacio de
soluciones es exactamente $S_n$.

\paragraph{Interpretación elegida.} El enunciado habla de una ``penalización
asociada'' al retraso. Se adoptó la lectura en la que la penalización es
\emph{proporcional al retraso} ($w_jT_j$). La alternativa, una multa fija si la
tarea llega tarde ($\sum w_jU_j$), es un problema distinto, más cercano a la
mochila. Sin fechas límite el problema sería $\sum w_jC_j$, que la regla de
Smith resuelve en tiempo polinomial (Sección~\ref{sec:wspt}); por eso las
fechas límite son esenciales para la dificultad.

\subsection{Ejemplo}\label{sec:ejemplo}

Considérense tres tareas:

\begin{center}
\begin{tabular}{cccc}
\toprule
Tarea & $p_j$ & $d_j$ & $w_j$\\
\midrule
A & 3 & 4 & 2\\
B & 2 & 2 & 5\\
C & 4 & 6 & 1\\
\bottomrule
\end{tabular}
\end{center}

Para $\pi=(A,B,C)$: $C_A=3$, $C_B=5$, $C_C=9$; los retrasos son $0$, $3$ y
$3$, y el costo es $2\cdot0+5\cdot3+1\cdot3=18$. La Tabla~\ref{tab:perm}
enumera las $3!=6$ permutaciones: la óptima es $(B,A,C)$ con costo $5$.
Poner primero la tarea más costosa de retrasar (B) reduce el costo a menos de
un tercio.

\begin{table}[h]
\centering
\caption{Todas las programaciones del ejemplo.}\label{tab:perm}
\begin{tabular}{lccc}
\toprule
$\pi$ & $(C_j)$ en orden & $(T_j)$ en orden & costo\\
\midrule
$A,B,C$ & $3,5,9$ & $0,3,3$ & 18\\
$A,C,B$ & $3,7,9$ & $0,1,7$ & 36\\
$\mathbf{B,A,C}$ & $2,5,9$ & $0,1,3$ & $\mathbf{5}$\\
$B,C,A$ & $2,6,9$ & $0,0,5$ & 10\\
$C,A,B$ & $4,7,9$ & $0,3,7$ & 41\\
$C,B,A$ & $4,6,9$ & $0,4,5$ & 30\\
\bottomrule
\end{tabular}
\end{table}

\subsection{Pertenencia a NP}

NP es una clase de \emph{problemas de decisión}, por lo que se considera la
versión:

\begin{quote}
\textsc{RetrasoPonderado}: dada una instancia y un entero $K$, ¿existe una
permutación $\pi$ con $\costo(\pi)\le K$?
\end{quote}

\begin{proposicion}
\textsc{RetrasoPonderado} $\in$ NP.
\end{proposicion}

\begin{proof}
El certificado es la permutación $\pi$, de tamaño $O(n\log n)$ bits, polinomial
en la entrada. El verificador comprueba que $\pi$ sea una permutación de $J$
(marcando cada índice en un arreglo, $O(n)$), recorre $\pi$ acumulando
$C_{\pi_k}=C_{\pi_{k-1}}+p_{\pi_k}$, suma $w_j\max(0,C_j-d_j)$ y compara el
total con $K$. Son $O(n)$ operaciones aritméticas sobre números de tamaño
polinomial en la entrada, así que la verificación es polinomial. Si la
respuesta es \emph{sí}, el $\pi$ correcto hace aceptar al verificador; si es
\emph{no}, ningún $\pi$ lo hace.
\end{proof}

El problema de \emph{optimización} no está en NP en sentido estricto, porque no
es de decisión, pero es tan difícil como la versión de decisión. Con un
algoritmo que resuelva la decisión se encuentra el óptimo por búsqueda binaria
sobre $K\in[0,\sum_j w_j\sum_i p_i]$, y conocido el costo óptimo se reconstruye
la permutación fijando las posiciones una a una.

\subsection{Complejidad computacional}

\begin{teorema}[Lawler 1977; Lenstra, Rinnooy Kan y Brucker 1977]
$1\,\|\,\sum w_jT_j$ es NP-difícil en sentido fuerte. En consecuencia,
\textsc{RetrasoPonderado} es NP-completo en sentido fuerte.
\end{teorema}

La demostración original reduce \textsc{3-Partición}, que es NP-completo en
sentido fuerte, a este problema. Que la dificultad sea \emph{fuerte} tiene una
consecuencia concreta para el diseño: a diferencia de la mochila o de
$1\,\|\,\sum w_jU_j$, \textbf{no existe un algoritmo pseudo-polinomial}, es
decir, polinomial en $n$ y en la magnitud de los números, a menos que
$\mathrm{P}=\mathrm{NP}$. Una tabla indexada por el tiempo, $O(n\sum p_j)$, no
puede funcionar en general. Por eso la programación dinámica de la
Sección~\ref{sec:dp} es exponencial en $n$: es lo esperable para un problema
fuertemente NP-difícil.

\paragraph{Tamaño del espacio de búsqueda.} Hay $n!$ programaciones. Por la
fórmula de Stirling, $n!\approx\sqrt{2\pi n}\,(n/e)^n$, que crece más rápido
que cualquier $c^n$. Para $n=20$ hay unos $2.4\times10^{18}$ órdenes; evaluar
cada uno en $O(n)$ con $10^9$ operaciones por segundo tomaría más de 1500
años.

%=======================================================================
\section{Solución por programación dinámica}\label{sec:dp}
%=======================================================================

\subsection{Razonamiento que lleva a la formulación}

El punto de partida es la fuerza bruta, que considera los $n!$ órdenes. La
pregunta clave es: \emph{¿qué información de un prefijo de la programación se
necesita realmente para continuar?}

Supóngase que ya se decidió qué tareas van primero, el conjunto $S$, aunque no
su orden interno. Se observan dos hechos:

\begin{enumerate}
  \item Cualquiera sea su orden interno, las tareas de $S$ ocupan exactamente
        el intervalo $[0,P(S)]$, con $P(S)=\sum_{j\in S}p_j$. Las tareas
        fuera de $S$ empiezan todas después de $P(S)$, así que su costo
        \textbf{no depende} del orden interno de $S$.
  \item La última tarea ejecutada dentro de $S$ termina en $P(S)$, también
        independientemente del orden de las demás.
\end{enumerate}

Por lo tanto, para el futuro solo importa \emph{qué} tareas ya se ejecutaron,
no \emph{en qué orden}. Esto reduce la información de estado de $n!$ órdenes a
$2^n$ subconjuntos. Es la misma idea del algoritmo de Held--Karp para el
problema del viajante.

\begin{lema}\label{lema:fin}
En cualquier orden de las tareas de $S$ ejecutado desde $t=0$ sin huecos, la
última tarea termina en $P(S)$.
\end{lema}

\begin{proof}
Su tiempo de finalización es la suma de las duraciones de todas las tareas
ejecutadas hasta ella inclusive, que son exactamente las de $S$.
\end{proof}

\subsection{Modelo matemático}

\begin{definicion}[Estado]
Para $S\subseteq J$, $f(S)$ es el costo mínimo de ejecutar \emph{exactamente}
las tareas de $S$ al inicio de la programación, en el mejor orden posible.
\end{definicion}

\paragraph{Recurrencia.} Se elige cuál tarea $j\in S$ va de última. Por el
Lema~\ref{lema:fin}, termina en $P(S)$, y lo anterior es un subproblema sobre
$S\setminus\{j\}$:
\begin{equation}\label{eq:rec}
\boxed{
\begin{aligned}
f(\varnothing)&=0,\\
f(S)&=\min_{j\in S}\Big[f(S\setminus\{j\})+w_j\max\big(0,\,P(S)-d_j\big)\Big].
\end{aligned}}
\end{equation}
La respuesta es $f(J)$. Para reconstruir el orden se guarda, para cada $S$, la
tarea $u(S)$ que alcanza el mínimo, y se retrocede desde $J$:
$u(J)$ es la última tarea, $u(J\setminus\{u(J)\})$ la penúltima, etc.

\begin{teorema}[Corrección]
Para todo $S\subseteq J$, $f(S)$ calculado con~\eqref{eq:rec} es igual al costo
mínimo de ordenar las tareas de $S$ desde $t=0$.
\end{teorema}

\begin{proof}
Por inducción sobre $|S|$. Si $S=\varnothing$ el costo es $0$. Sea
$|S|\ge1$ y llámese $\mathrm{OPT}(S)$ al costo mínimo verdadero.

($\ge$) Sea $\sigma$ un orden óptimo de $S$ y $j$ su última tarea. Por el
Lema~\ref{lema:fin}, $j$ cuesta $w_j\max(0,P(S)-d_j)$. Las demás forman un orden
de $S\setminus\{j\}$ que empieza en $t=0$, de costo
$\ge\mathrm{OPT}(S\setminus\{j\})=f(S\setminus\{j\})$ por hipótesis
inductiva. Entonces $\mathrm{OPT}(S)\ge f(S\setminus\{j\})+w_j\max(0,P(S)-d_j)\ge f(S)$.

($\le$) Sea $j$ la tarea que minimiza~\eqref{eq:rec}. Tomando un orden óptimo
de $S\setminus\{j\}$, cuyo costo es $f(S\setminus\{j\})$ por hipótesis
inductiva, y poniendo $j$ al final, se obtiene un orden de $S$ con costo
exactamente $f(S)$. Así $\mathrm{OPT}(S)\le f(S)$.
\end{proof}

La demostración evidencia las dos propiedades que justifican la programación
dinámica: \textbf{subestructura óptima}, porque el prefijo de un orden óptimo
es óptimo para su conjunto, y \textbf{subproblemas superpuestos}, porque el
mismo $f(S)$ se usa desde todos los $S\cup\{j\}$, hasta $n-|S|$ veces.

\subsection{Ejemplo: tabla completa}

La Tabla~\ref{tab:dp} aplica~\eqref{eq:rec} al ejemplo de la
Sección~\ref{sec:ejemplo}. Cada candidato es
$f(S\setminus\{j\})+w_j\max(0,P(S)-d_j)$ con $j$ como última tarea.

\begin{table}[h]
\centering
\caption{Programación dinámica sobre el ejemplo A, B, C.}\label{tab:dp}
\footnotesize\setlength{\tabcolsep}{4pt}
\begin{tabular}{lcll}
\toprule
$S$ & $P(S)$ & candidatos ($j$: valor) & $f(S)$\\
\midrule
$\varnothing$ & 0 & -- & 0\\
$\{A\}$ & 3 & A: $0+2\cdot0=0$ & 0\\
$\{B\}$ & 2 & B: $0+5\cdot0=0$ & 0\\
$\{C\}$ & 4 & C: $0+1\cdot0=0$ & 0\\
$\{A,B\}$ & 5 & A: $0+2\cdot1=\mathbf{2}$;\; B: $0+5\cdot3=15$ & 2\\
$\{A,C\}$ & 7 & A: $0+2\cdot3=6$;\; C: $0+1\cdot1=\mathbf{1}$ & 1\\
$\{B,C\}$ & 6 & B: $0+5\cdot4=20$;\; C: $0+1\cdot0=\mathbf{0}$ & 0\\
$\{A,B,C\}$ & 9 & A: $0+10=10$;\; B: $1+35=36$; & \\
 & & C: $2+3=\mathbf{5}$ & 5\\
\bottomrule
\end{tabular}
\end{table}

Reconstrucción:
\[
u(\{A,B,C\})=C,\quad u(\{A,B\})=A,\quad u(\{B\})=B.
\]
El orden es $(B,A,C)$ con costo $5$, igual al de la Tabla~\ref{tab:perm}. Se
evaluaron $12$ candidatos en lugar de las $6\cdot3=18$ operaciones de la fuerza
bruta; la diferencia crece de forma exponencial con $n$.

\subsection{Algoritmo}

Los subconjuntos se representan como máscaras de bits: el bit $j$ de un entero
$S$ vale $1$ si la tarea $j$ está en $S$. Si $S'\subset S$, entonces $S'<S$
como entero, así que basta recorrer las máscaras en orden creciente para que
todo $f(S\setminus\{j\})$ esté calculado antes de usarlo.

\begin{algorithm}[h]
\caption{Programación dinámica sobre subconjuntos}\label{alg:dp}
\begin{algorithmic}[1]
\Require tareas $(p_j,d_j,w_j)$, $j=0,\dots,n-1$
\Ensure orden óptimo y su costo
\State $f[0]\gets0$
\For{$S\gets1$ \textbf{hasta} $2^n-1$}
  \State $P\gets\sum_{j\in S}p_j$; \quad $f[S]\gets\infty$
  \For{cada $j\in S$}
    \State $c\gets f[S\setminus\{j\}]+w_j\max(0,P-d_j)$
    \If{$c<f[S]$} $f[S]\gets c$;\; $u[S]\gets j$ \EndIf
  \EndFor
\EndFor
\State $S\gets2^n-1$;\; $\pi\gets()$
\While{$S\ne0$}
  \State anteponer $u[S]$ a $\pi$;\; $S\gets S\setminus\{u[S]\}$
\EndWhile
\State \Return $(\pi,\,f[2^n-1])$
\end{algorithmic}
\end{algorithm}

\subsection{Complejidad}

\paragraph{Tiempo.} Hay $2^n$ estados; para cada uno se calcula $P(S)$ en
$O(n)$ y se prueban hasta $n$ candidatos en $O(1)$ cada uno. Total
\[
T(n)=\Theta(2^n\cdot n).
\]
Más precisamente, el número de candidatos es
$\sum_{S}|S|=\sum_{k=0}^{n}k\binom{n}{k}=n\,2^{n-1}$. La reconstrucción cuesta
$O(n)$.

\paragraph{Espacio.} Se guardan $f[S]$ (entero de 64 bits) y $u[S]$ (entero de
8 bits) para cada $S$: $\Theta(2^n)$, unos $9\cdot2^n$ bytes. Para $n=24$ son
cerca de 150\,MB, que se tomó como límite práctico de la implementación.

\paragraph{Comparación con la fuerza bruta.} La fuerza bruta cuesta
$\Theta(n\cdot n!)$. El cociente
\[
\frac{n\cdot n!}{n\cdot2^n}=\frac{n!}{2^n}=\prod_{k=1}^{n}\frac{k}{2}
\]
crece sin cota: para $n=10$ es $\approx3.5\times10^3$ y para $n=20$ es
$\approx2.3\times10^{12}$. La programación dinámica sigue siendo exponencial,
como se espera de un problema fuertemente NP-difícil, pero mueve la frontera
de lo resoluble de $n\approx11$ a $n\approx25$.

%=======================================================================
\section{Solución mediante enfoques voraces}
%=======================================================================

Un algoritmo voraz construye la programación eligiendo en cada paso la tarea
que \emph{parece} mejor según un criterio local, sin reconsiderar decisiones.
Se implementaron tres criterios. Los dos primeros son reglas clásicas que
ignoran uno de los parámetros; el tercero, WMDD, combina ambos y es el voraz
principal de este trabajo. En todos los casos los empates se resuelven a favor
de la tarea de menor índice.

\subsection{EDD: fecha límite más cercana}

\textbf{Criterio:} ordenar por $d_j$ creciente. \textbf{Complejidad:}
$O(n\log n)$.

\textbf{Justificación:} atender primero lo que vence primero. Por el teorema de
Jackson (1955), EDD minimiza el retraso \emph{máximo}; en particular, si
existe una programación en la que ninguna tarea se retrasa, EDD la encuentra.

\textbf{Limitación:} ignora por completo los pesos $w_j$ y las duraciones, así
que una tarea barata con fecha cercana puede desplazar a una muy costosa.

\subsection{WSPT: regla de Smith}\label{sec:wspt}

\textbf{Criterio:} ordenar por $w_j/p_j$ decreciente, es decir, mayor
penalización por unidad de tiempo ocupado. \textbf{Complejidad:}
$O(n\log n)$.

\begin{proposicion}[Intercambio adyacente]\label{prop:intercambio}
Sean $i,j$ tareas consecutivas que empiezan en el instante $t$, ambas tardías
en cualquiera de sus dos órdenes, lo que ocurre si $d_i\le t+p_i$ y
$d_j\le t+p_j$. Entonces conviene poner $i$ antes que $j$ si y solo si
$w_i/p_i\ge w_j/p_j$.
\end{proposicion}

\begin{proof}
Si ambas son tardías, $T=C-d$. Con $i$ primero, el costo del par es
$w_i(t+p_i-d_i)+w_j(t+p_i+p_j-d_j)$; con $j$ primero,
$w_j(t+p_j-d_j)+w_i(t+p_j+p_i-d_i)$. La resta (primero menos segundo) es
$w_jp_i-w_ip_j$, que es $\le0$ si y solo si $w_i/p_i\ge w_j/p_j$. El resto de
la programación no cambia porque el par termina en $t+p_i+p_j$ en ambos casos.
\end{proof}

\begin{observacion}
Si todas las fechas límite son $0$, todas las tareas son tardías y el objetivo
se vuelve $\sum w_jC_j-\text{cte}$. Aplicando la
Proposición~\ref{prop:intercambio} a cualquier par adyacente fuera del orden
WSPT se demuestra que WSPT es \emph{óptimo} en ese caso (Smith, 1956). Por eso
el problema sin fechas límite está en P.
\end{observacion}

\textbf{Limitación:} ignora las fechas límite; puede adelantar una tarea con
mucha holgura y atrasar otra que estaba por vencer.

\subsection{WMDD: fecha límite modificada ponderada}

\textbf{Criterio:} en cada paso, con $t$ igual al tiempo ya ocupado, elegir la
tarea pendiente que minimice
\[
\rho_j(t)=\frac{\max\big(p_j,\;d_j-t\big)}{w_j}.
\]
Las tareas con $w_j=0$ no generan costo, tienen prioridad $+\infty$ y van al
final.

\textbf{Justificación del diseño.} $\rho_j$ cambia de régimen según la
holgura $d_j-t$ de cada tarea:
\begin{itemize}
  \item Si la tarea ya no puede terminar a tiempo ($d_j-t\le p_j$), entonces
        $\rho_j=p_j/w_j$: es WSPT, que la Proposición~\ref{prop:intercambio}
        justifica entre tareas tardías.
  \item Si aún tiene holgura ($d_j-t>p_j$), entonces $\rho_j=(d_j-t)/w_j$: es
        EDD ponderado, donde se adelanta lo que vence antes, ajustado por lo
        que cuesta retrasarlo.
\end{itemize}
El criterio es \emph{dinámico}: como depende de $t$, debe recalcularse en cada
paso. Por eso cuesta $O(n^2)$, con $n$ pasos y $n$ candidatos por paso. No se
usan números de punto flotante: $a/w<a'/w'$ se compara como $a\,w'<a'\,w$, lo
que con los límites de entrada ($a\le10^9$, $w\le10^6$) cabe en 64 bits.

\begin{algorithm}[h]
\caption{Voraz WMDD}\label{alg:wmdd}
\begin{algorithmic}[1]
\Require tareas $(p_j,d_j,w_j)$, $j=1,\dots,n$
\Ensure un orden $\pi$ (no necesariamente óptimo)
\State $R\gets\{1,\dots,n\}$;\; $t\gets0$;\; $\pi\gets()$
\While{$R\ne\varnothing$}
  \State $k\gets\arg\min_{j\in R}\;\max(p_j,\,d_j-t)/w_j$
  \Statex \hspace{\algorithmicindent}(empates: menor índice; $w_j=0$ cuenta como $+\infty$)
  \State agregar $k$ al final de $\pi$
  \State $t\gets t+p_k$;\; $R\gets R\setminus\{k\}$
\EndWhile
\State \Return $\pi$
\end{algorithmic}
\end{algorithm}

La línea 3 es la única decisión del algoritmo y nunca se revisa después. Esa
es la diferencia esencial con el Algoritmo~\ref{alg:dp}, que para cada conjunto
compara \emph{todas} las tareas que podrían ir de últimas.

\subsection{¿Garantizan optimalidad?}\label{sec:optimalidad}

\textbf{No.} Ningún voraz polinomial puede garantizarla: si alguno lo hiciera
se tendría un algoritmo polinomial para un problema NP-difícil, y por tanto
P\,=\,NP. Basta un contraejemplo concreto.

\paragraph{Contraejemplo.} Sean las tareas
\begin{center}
\begin{tabular}{ccccc}
\toprule
Tarea & 1 & 2 & 3 & 4\\
\midrule
$p_j$ & 4 & 5 & 6 & 4\\
$d_j$ & 9 & 7 & 6 & 5\\
$w_j$ & 5 & 3 & 3 & 2\\
\bottomrule
\end{tabular}
\end{center}
El óptimo, verificado por fuerza bruta y por programación dinámica, es
$(2,1,3,4)$ con costo $55$: las tareas 2 y 1 terminan en $5$ y $9$, a tiempo, y
3 y 4 se retrasan $9$ y $14$ unidades ($3\cdot9+2\cdot14=55$).

WMDD en $t=0$ calcula $\rho=(9/5,\,7/3,\,6/3,\,5/2)=(1.8,\,2.33,\,2,\,2.5)$ y
elige la tarea 1, por su peso alto. Pero la tarea 1 tenía holgura de sobra:
podía ir segunda y aun así terminar en $9=d_1$. Elegirla primero retrasa la 2
en $2$ unidades. El resultado es $(1,2,3,4)$ con costo $61$, una brecha de
$10.9$\,\%. WSPT produce el mismo orden (61) y EDD obtiene $(4,3,2,1)$ con
costo $86$ ($+56.4$\,\%).

El error tiene una causa clara: el voraz mira solo el instante actual y no
\emph{prevé} que retrasar una tarea con holgura no cuesta nada. La
programación dinámica sí lo considera, porque evalúa todos los conjuntos
posibles de tareas previas.

\subsection{Comparación conceptual con la programación dinámica}

\begin{table}[h]
\centering
\caption{Programación dinámica frente a voraz.}\label{tab:comp}
\small
\setlength{\tabcolsep}{3pt}\begin{tabular}{>{\raggedright\arraybackslash}p{1.8cm}>{\raggedright\arraybackslash}p{2.8cm}>{\raggedright\arraybackslash}p{2.8cm}}
\toprule
 & Programación dinámica & Voraz (WMDD)\\
\midrule
Decisiones & considera todas las tareas que pueden ir de últimas en cada conjunto & una elección local e irrevocable por paso\\
Optimalidad & garantizada (teorema de corrección) & no garantizada (contraejemplo)\\
Tiempo & $\Theta(2^n n)$ & $\Theta(n^2)$\\
Espacio & $\Theta(2^n)$ & $\Theta(n)$\\
Tamaño práctico & $n\le24$ & miles de tareas\\
Principio & subestructura óptima y subproblemas superpuestos & criterio de intercambio local (Prop.~\ref{prop:intercambio})\\
\bottomrule
\end{tabular}
\end{table}

Ambos enfoques parten de observaciones estructurales distintas. La
programación dinámica aprovecha que el costo futuro depende solo del
\emph{conjunto} ya programado, y por eso puede explorar todas las opciones sin
repetir trabajo. El voraz aprovecha que, entre dos tareas tardías adyacentes,
el orden correcto lo decide una razón local ($w/p$). Esa regla es exacta para
pares de tareas tardías, pero no captura la interacción global entre holguras,
que es justamente lo que hace difícil al problema.

%=======================================================================
\section{Implementación}
%=======================================================================

Ambas soluciones se implementaron en C++17, sin bibliotecas externas, y se
compilan con \texttt{g++ -O2} mediante un \texttt{Makefile}. El código está
organizado en módulos pequeños con una interfaz común: cada algoritmo recibe el
vector de tareas y devuelve un \texttt{Resultado} con el orden (identificadores
$1..n$) y su costo.

\begin{center}
\small
\begin{tabular}{ll}
\toprule
Archivo & Responsabilidad\\
\midrule
\texttt{entrada.cpp} & lectura y validación de instancias\\
\texttt{costo.cpp} & costo y detalle de un orden\\
\texttt{dp.cpp} & programación dinámica\\
\texttt{fuerza\_bruta.cpp} & verificación por enumeración\\
\texttt{voraz.cpp} & EDD, WSPT y WMDD\\
\texttt{generacion.cpp} & instancias aleatorias\\
\texttt{main.cpp} & programa \texttt{tareas}\\
\texttt{experimentos.cpp} & mediciones en CSV\\
\bottomrule
\end{tabular}
\end{center}

\paragraph{Decisiones de implementación relevantes.}
\begin{itemize}
  \item \textbf{Máscaras de 32 bits} para los subconjuntos y recorrido en
        orden creciente, lo que evita la recursión y la memoización con mapas.
        $P(S)$ se calcula en el mismo ciclo que los candidatos, sin un
        arreglo adicional.
  \item \textbf{Arreglo de predecesores de 8 bits} ($u[S]\in[0,23]$), que
        reduce la memoria de $16$ a $9$ bytes por estado.
  \item \textbf{Aritmética entera exacta}: costos en \texttt{long long} y
        comparaciones de cocientes por producto cruzado. Los límites de
        entrada ($n\le1000$, $p,w\le10^6$, $d\le10^9$) garantizan que no haya
        desbordamiento; el peor caso, $\approx5\times10^{17}$, se verifica en
        las pruebas.
  \item \textbf{Validación de la entrada} con mensajes que indican la línea
        del error. Se toleran comentarios, líneas vacías, fines de línea de
        Windows y la marca BOM, para poder usar archivos preparados en otro
        sistema durante la sustentación.
  \item \textbf{Límites explícitos}: la fuerza bruta rechaza $n>10$ y la
        programación dinámica $n>24$ con un mensaje claro, en vez de
        agotar la memoria.
\end{itemize}

\paragraph{Uso.} \texttt{./tareas archivo.txt [--algoritmo dp|fb|edd|wspt|wmdd|todos] [--detalle]}
muestra para cada algoritmo el orden, el costo y el tiempo y, con
\texttt{--detalle}, la tabla $C_j,T_j,w_jT_j$ de cada tarea. Al final reporta
la brecha de cada voraz frente al óptimo. El archivo de entrada contiene $n$ y
luego una línea $p_j\;d_j\;w_j$ por tarea. \texttt{./generador n semilla} crea
instancias reproducibles y \texttt{./experimentos} produce los datos de la
Sección~\ref{sec:pruebas}.

%=======================================================================
\section{Pruebas y resultados}\label{sec:pruebas}
%=======================================================================

\subsection{Pruebas de corrección}

El conjunto automático (\texttt{make test}, 2558 comprobaciones) incluye:
\begin{itemize}
  \item Costos calculados a mano (ejemplo A, B, C y su detalle).
  \item \textbf{Programación dinámica igual a fuerza bruta} en 300 instancias
        aleatorias con $n\le8$: el costo coincide en todas.
  \item Todo orden devuelto es una permutación de $1..n$ y su costo
        recalculado coincide con el reportado.
  \item Ningún voraz obtiene un costo menor que el de la programación
        dinámica, y existe al menos una instancia donde WMDD no es óptimo.
  \item Casos borde: $n=0$, pesos nulos, valores máximos, límites de $n$ y
        archivos mal formados, que deben terminar con código de error 1.
\end{itemize}

\subsection{Casos diseñados por el equipo}

Además de las pruebas automáticas, se diseñaron cuatro instancias
(carpeta \texttt{ejemplos/}) para ilustrar comportamientos concretos. La
Tabla~\ref{tab:casos} muestra el costo obtenido por cada algoritmo al ejecutar
\texttt{./tareas} sobre cada archivo.

\begin{table}[h]
\centering
\caption{Costo de cada algoritmo en los casos del equipo (FB = fuerza bruta,
PD = programación dinámica; en negrita, los que alcanzan el óptimo).}\label{tab:casos}
\small\setlength{\tabcolsep}{3pt}
\begin{tabular}{lrrrrrr}
\toprule
Caso & $n$ & FB & PD & EDD & WSPT & WMDD\\
\midrule
\texttt{ejemplo\_basico}   & 3  & \textbf{5}  & \textbf{5}  & \textbf{5}  & \textbf{5}  & \textbf{5}\\
\texttt{todas\_a\_tiempo}  & 4  & \textbf{0}  & \textbf{0}  & \textbf{0}  & \textbf{0}  & \textbf{0}\\
\texttt{voraz\_no\_optimo} & 4  & \textbf{55} & \textbf{55} & 86 & 61 & 61\\
\texttt{mediano\_15}       & 15 & --          & \textbf{798} & 1495 & 1028 & \textbf{798}\\
\bottomrule
\end{tabular}
\end{table}

\begin{itemize}
  \item \textbf{\texttt{ejemplo\_basico}}: el ejemplo de la
        Sección~\ref{sec:ejemplo}. Todos los algoritmos encuentran el orden
        $(B,A,C)$: con pocas tareas y criterios alineados (B vence primero y
        es la más costosa de retrasar), los voraces aciertan.
  \item \textbf{\texttt{todas\_a\_tiempo}}: fechas límite holgadas, así que
        cualquier orden cuesta $0$. Los algoritmos devuelven órdenes distintos,
        todos óptimos: el óptimo no tiene por qué ser único.
  \item \textbf{\texttt{voraz\_no\_optimo}}: el contraejemplo de la
        Sección~\ref{sec:optimalidad}. La fuerza bruta devuelve $(2,1,3,4)$ y
        la programación dinámica $(2,1,4,3)$, ambos de costo $55$: dos óptimos
        distintos que cada algoritmo desempata de forma diferente.
  \item \textbf{\texttt{mediano\_15}}: instancia aleatoria
        (\texttt{./generador 15 2024}) fuera del alcance de la fuerza bruta
        ($15!\approx1.3\times10^{12}$ órdenes). La programación dinámica la
        resuelve en milisegundos y WMDD también llega al óptimo, mientras que
        EDD casi duplica el costo ($+87$\,\%).
\end{itemize}

También se comprobó el manejo de entradas inválidas: un archivo con menos
tareas de las declaradas, valores negativos o no numéricos, y la petición de
programación dinámica con $n=25$ terminan con un mensaje que indica la causa
y código de salida $1$.

\subsection{Diseño experimental}

Se generaron instancias con $p_j,w_j\sim U\{1,\dots,10\}$ y
$d_j\sim U\{0,\dots,\lfloor\alpha\sum p_i\rfloor\}$, donde $\alpha$ controla
la \emph{holgura} de las fechas límite: alta ($\alpha=1$), media ($0.6$) y baja
($0.3$, la mayoría de las tareas termina tarde). Para cada $n\in\{4,\dots,24\}$
se usaron 5 semillas y las 3 holguras, es decir $21\cdot5\cdot3=315$
instancias. La fuerza bruta solo se ejecutó con $n\le10$. Los tiempos se
midieron con \texttt{std::chrono::steady\_clock}, en un solo hilo, en un
computador con Linux.

\subsection{Tiempo de ejecución}

\begin{figure}[h]
\centering
\begin{tikzpicture}
\begin{axis}[
  width=\columnwidth, height=6.2cm,
  ymode=log, xlabel={$n$ (número de tareas)}, ylabel={tiempo promedio (ms)},
  xmin=3.5, xmax=24.5, legend pos=south east, legend style={font=\small},
  grid=major, grid style={gray!25}, tick label style={font=\small}]
\addplot[mark=*,mark size=1.5pt,blue!70!black,thick] coordinates {
 (4,0.0004)(5,0.0009)(6,0.0024)(7,0.0057)(8,0.0149)(9,0.0331)(10,0.0701)
 (11,0.1456)(12,0.3043)(13,0.6292)(14,1.343)(15,2.8149)(16,5.8534)
 (17,12.2872)(18,26.1342)(19,53.6211)(20,111.8425)(21,229.9854)
 (22,474.7733)(23,983.2248)(24,1987.3369)};
\addlegendentry{Prog. dinámica}
\addplot[mark=square*,mark size=1.5pt,red!75!black,thick] coordinates {
 (4,0.0012)(5,0.007)(6,0.0401)(7,0.2885)(8,2.3375)(9,24.9574)(10,268.3976)};
\addlegendentry{Fuerza bruta}
\addplot[mark=triangle*,mark size=1.8pt,green!50!black,thick] coordinates {
 (4,0.0002)(6,0.0003)(8,0.0004)(10,0.0008)(12,0.0007)(14,0.0009)
 (16,0.0011)(18,0.0028)(20,0.0019)(22,0.0018)(24,0.0023)};
\addlegendentry{Voraz WMDD}
\end{axis}
\end{tikzpicture}
\caption{Tiempo promedio por instancia (escala logarítmica).}\label{fig:tiempo}
\end{figure}

\begin{table}[h]
\centering
\caption{Tiempo promedio (ms) y razón de crecimiento al aumentar $n$ en 1.}\label{tab:tiempo}
\small
\begin{tabular}{rrrrr}
\toprule
$n$ & PD & razón & FB & razón\\
\midrule
7 & 0.0057 & 2.4 & 0.289 & 7.2\\
8 & 0.0149 & 2.6 & 2.34 & 8.1\\
9 & 0.0331 & 2.2 & 25.0 & 10.7\\
10 & 0.0701 & 2.1 & 268.4 & 10.8\\
\midrule
20 & 111.8 & 2.09 & -- & --\\
21 & 230.0 & 2.06 & -- & --\\
22 & 474.8 & 2.06 & -- & --\\
23 & 983.2 & 2.07 & -- & --\\
24 & 1987.3 & 2.02 & -- & --\\
\bottomrule
\end{tabular}
\end{table}

La Figura~\ref{fig:tiempo} y la Tabla~\ref{tab:tiempo} confirman el análisis
asintótico:
\begin{itemize}
  \item En escala logarítmica, la programación dinámica es casi una recta: el
        tiempo \textbf{se multiplica por $\approx2.05$} con cada tarea
        adicional. El modelo $\Theta(2^n n)$ predice
        $2(n+1)/n\approx2.09$ para $n\approx22$.
  \item La fuerza bruta se multiplica por $\approx n+1$, como predice
        $\Theta(n\cdot n!)$, y su curva se curva hacia arriba. Con $n=10$ ya es
        unas 3800 veces más lenta que la programación dinámica; extrapolando,
        con $n=14$ tardaría horas.
  \item Los voraces se mantienen por debajo de $0.003$\,ms en todo el rango;
        a esta escala el tiempo está dominado por costos fijos (reserva de
        memoria, medición), no por el $n^2$.
  \item Con $n=24$ la programación dinámica tarda 2\,s y usa cerca de
        150\,MB. Cada tarea extra duplica ambos, así que $n\approx28$ ya
        requeriría más de 2\,GB.
\end{itemize}

\subsection{Calidad de los voraces}

Se define la brecha de un voraz como
\[
\text{brecha}=100\cdot\frac{\costo_{\text{voraz}}-\costo_{\text{óptimo}}}{\costo_{\text{óptimo}}}.
\]

\begin{table}[h]
\centering
\caption{Brecha (\%) de cada voraz frente al óptimo, 105 instancias por holgura.}\label{tab:brecha}
\small\setlength{\tabcolsep}{4pt}
\begin{tabular}{llrrrr}
\toprule
Voraz & Holgura & media & mediana & máx. & óptimos\\
\midrule
EDD  & alta  & 111.1 & 84.5 & 450.4 & 10\\
     & media & 102.0 & 96.3 & 255.1 & 1\\
     & baja  & 72.2  & 73.0 & 170.9 & 1\\
\midrule
WSPT & alta  & 495.7 & 227.1 & 10500 & 3\\
     & media & 30.2  & 29.1  & 84.2  & 7\\
     & baja  & 4.2   & 4.0   & 10.8  & 20\\
\midrule
WMDD & alta  & 115.9 & 17.6 & 7500  & 12\\
     & media & \textbf{2.0} & \textbf{1.2} & 13.1 & 27\\
     & baja  & \textbf{0.6} & \textbf{0.2} & 4.8  & 44\\
\bottomrule
\end{tabular}
\end{table}

De la Tabla~\ref{tab:brecha} se concluye:
\begin{itemize}
  \item \textbf{WMDD domina} a los otros dos voraces en todas las holguras. Con
        holgura media y baja su brecha mediana es de $1.2$\,\% y $0.2$\,\%, y
        alcanza el óptimo exacto en 71 de 210 instancias.
  \item \textbf{WSPT mejora a medida que la holgura baja}: de $227$\,\% a
        $4$\,\% de mediana. Es lo que predice la teoría, porque cuando casi
        todas las tareas son tardías la Proposición~\ref{prop:intercambio}
        se aplica a casi todos los pares.
  \item \textbf{EDD es el peor} cuando hay retrasos, porque ignora los pesos.
        Su brecha mediana nunca baja de $73$\,\%.
  \item \textbf{Con holgura alta los porcentajes engañan}: el costo óptimo
        suele ser muy pequeño (pocas tareas tardías), y una diferencia
        absoluta de pocas unidades produce brechas de miles por ciento (el
        máximo de 7500\,\% de WMDD es una instancia con óptimo $1$ donde el voraz obtiene $76$). La
        mediana es una medida más representativa que la media en este
        régimen.
\end{itemize}

\subsection{Comportamiento en función del tamaño de la entrada}

\begin{itemize}
  \item La \textbf{programación dinámica} es exacta para cualquier $n$, pero su
        costo crece exponencialmente. El límite es la memoria más que el
        tiempo.
  \item La \textbf{brecha del voraz} no muestra una tendencia monótona con
        $n$: la media de WMDD oscila entre $4$\,\% y $30$\,\% para
        $n\ge9$, salvo un pico en $n=15$. Depende mucho más de la estructura de la instancia (la
        holgura) que de su tamaño. Los picos aislados, como $n=8$ o $n=15$,
        vienen de instancias de holgura alta con óptimo cercano a cero.
  \item Sí aumenta con $n$ la \emph{probabilidad} de no dar el óptimo: con
        $n=4$ WMDD es óptimo en 13 de 15 instancias, mientras que con
        $n\ge20$ solo lo fue en 5 de 75, porque hay más pares de tareas donde una
        decisión local puede equivocarse.
\end{itemize}

%=======================================================================
\section{Análisis del proceso}
%=======================================================================

\paragraph{Cómo se llegó a cada modelo.} El primer paso fue precisar el
problema. El enunciado admitía dos lecturas (penalización proporcional al
retraso o multa fija), y se eligió la proporcional por ser la más literal. Se
comprobó además que sin fechas límite el problema sería polinomial, lo que
justificó incluirlas en el modelo.

Para la programación dinámica se partió de la fuerza bruta y se buscó qué parte
del prefijo de un orden afecta al resto. Al notar que la última tarea de un
conjunto siempre termina en $P(S)$ (Lema~\ref{lema:fin}), el estado natural
resultó ser el conjunto y no el orden. Se descartó un estado indexado por el
tiempo, $f(k,t)$, porque no permite saber qué tareas faltan; esa es la razón
de fondo por la que no hay solución pseudo-polinomial.

Para el voraz se analizó primero por qué fallan las reglas simples: EDD
ignora $w$ y WSPT ignora $d$. El argumento de intercambio
(Proposición~\ref{prop:intercambio}) mostró que WSPT es correcto entre tareas
tardías, y eso llevó a un criterio que se comporta como WSPT para las tareas
tardías y como EDD ponderado para las que aún tienen holgura: WMDD.

\paragraph{Dificultades encontradas.}
\begin{itemize}
  \item \textbf{Ambigüedad del enunciado}, resuelta con el análisis de las
        variantes y de su complejidad.
  \item \textbf{Memoria de la programación dinámica}: guardar costo,
        predecesor y $P(S)$ en arreglos de 64 bits costaría 24 bytes por
        estado. Por eso se usó un predecesor de 8 bits y se calculó $P(S)$
        dentro del ciclo: 9 bytes por estado, lo que permite llegar a $n=24$.
  \item \textbf{Archivos de entrada de otros sistemas}: un archivo guardado
        en Windows como UTF-8 con BOM era rechazado con un error confuso. Se
        detectó en la revisión del código y se corrigió con una prueba que lo
        cubre.
  \item \textbf{Tareas con $w_j=0$}: su prioridad en WMDD implicaría dividir
        por cero. Se definió prioridad infinita (no cuesta retrasarlas) y se
        evitó la división con productos cruzados.
  \item \textbf{Verificar la optimalidad}: no basta con probar algunos ejemplos
        a mano. La fuerza bruta se usó como oráculo sobre cientos de instancias
        aleatorias.
  \item \textbf{Interpretar la brecha}: la media en porcentaje resultó
        engañosa con óptimos pequeños, y se complementó con la mediana.
\end{itemize}

\paragraph{Decisiones de diseño.} C++ por su control de memoria y su
eficiencia en los ciclos de la programación dinámica; máscaras de bits en lugar
de conjuntos explícitos; los mismos módulos para el programa, las pruebas y
los experimentos; instancias generadas con semilla fija para que todos los
resultados sean reproducibles.

%=======================================================================
\section{Conclusiones}
%=======================================================================

\begin{enumerate}
  \item \textbf{La programación dinámica no elimina la dificultad, la
        reorganiza.} Pasar de $n!$ órdenes a $2^n$ conjuntos es una mejora
        enorme, del orden de $10^{12}$ para $n=20$, pero sigue siendo
        exponencial. Esto es coherente con que el problema sea fuertemente
        NP-difícil: no cabe esperar una tabla polinomial ni pseudo-polinomial.
        La ganancia proviene de identificar la información \emph{mínima}
        suficiente sobre el pasado, el conjunto y no el orden.
  \item \textbf{La frontera práctica la fija la memoria.} El tiempo se duplica
        y la memoria también con cada tarea; con $n=24$ la memoria se acerca
        al límite razonable antes que el tiempo. En la práctica, el análisis de
        espacio es tan importante como el de tiempo.
  \item \textbf{Un voraz es tan bueno como la estructura que explota.} WSPT es
        exacto cuando todas las tareas son tardías y empeora cuando no lo son;
        WMDD funciona bien porque su criterio cambia de régimen siguiendo esa
        misma estructura. La calidad de un voraz no es una propiedad fija: depende
        de la distribución de las instancias.
  \item \textbf{La falta de optimalidad es estructural, no un defecto de
        implementación.} El contraejemplo muestra el mecanismo: una decisión
        local no puede prever que retrasar una tarea con holgura es gratis.
        Un voraz que siempre acertara probaría P\,=\,NP.
  \item \textbf{Los dos enfoques se complementan.} La programación dinámica
        sirve como referencia exacta para medir los voraces, y estos resuelven
        en microsegundos instancias que la programación dinámica no puede
        abordar. Una extensión natural sería combinarlos: usar WMDD como cota
        superior inicial en un \emph{branch and bound}, o como punto de partida
        de una búsqueda local.
  \item \textbf{Medir bien es parte del análisis.} La brecha porcentual con
        óptimos cercanos a cero produjo cifras engañosas. Reportar medianas y
        separar los resultados por tipo de instancia cambió las conclusiones
        respecto a mirar solo promedios globales.
\end{enumerate}

\begin{thebibliography}{9}
\small
\bibitem{lawler} E.~L. Lawler. A ``pseudopolynomial'' algorithm for sequencing
jobs to minimize total tardiness. \emph{Annals of Discrete Mathematics}, 1:331--342, 1977.
\bibitem{lenstra} J.~K. Lenstra, A.~H.~G. Rinnooy Kan y P. Brucker. Complexity of
machine scheduling problems. \emph{Annals of Discrete Mathematics}, 1:343--362, 1977.
\bibitem{smith} W.~E. Smith. Various optimizers for single-stage production.
\emph{Naval Research Logistics Quarterly}, 3:59--66, 1956.
\bibitem{jackson} J.~R. Jackson. Scheduling a production line to minimize maximum
tardiness. Research Report 43, UCLA, 1955.
\bibitem{heldkarp} M. Held y R.~M. Karp. A dynamic programming approach to
sequencing problems. \emph{Journal of SIAM}, 10(1):196--210, 1962.
\bibitem{pinedo} M.~L. Pinedo. \emph{Scheduling: Theory, Algorithms, and Systems}.
Springer, 5.ª ed., 2016.
\bibitem{clrs} T.~H. Cormen, C.~E. Leiserson, R.~L. Rivest y C. Stein.
\emph{Introduction to Algorithms}. MIT Press, 4.ª ed., 2022.
\end{thebibliography}

\end{document}
